\subsection{FUNDAMENTAL DOMAIN OF W IN THE UNION OF THE CHAMBERS}

We retain the notation of no. 4. For \(s \in S\) , denote by \(H_s\) the hyperplane of \(E^*\) orthogonal to \(e_s\) , by \(\overline{A_s}\) the set of \(x^* \in E^*\) such that \(\langle x^*, e_s \rangle \geqslant 0\) and by \(\overline{C}\) the intersection of the \(\overline{A_s}\) for \(s \in S\) . For the weak topology \(\sigma(E^*, E)\) defined by the duality between \(E^*\) and \(E\) (Topological Vector Spaces, Chap. II, §6, no. 2), the \(\overline{A_s}\) are closed half-spaces and \(\overline{C}\) is a closed convex cone. Moreover, \(\overline{C}\) is the closure of C; indeed, if \(x^* \in \overline{C}\) and \(y^* \in C\) , then \(x^* + ty^* \in C\) for every real number \(t > 0\) and \(x^* = \lim_{t \to 0} (x^* + ty^*)\) .

For \(X\subset S\) ,put

\[
C _ {X} = \left(\bigcap_ {s \in X} H _ {s}\right) \cap \left(\bigcap_ {s \in S - X} A _ {s}\right).
\]

We have \(C_X \subset \overline{C}\) , \(C_{\varnothing} = C\) and \(C_S = \{0\}\) . The sets \(C_X\) , for \(X \in B(S)\) , form a partition of \(\overline{C}\) .

On the other hand, recall (Chap. IV, §1, no. 8) that \(\mathbf{W}_{\mathbf{X}}\) denotes the subgroup of \(\mathbf{W}\) generated by \(\mathbf{X}\) . Clearly \(w(x^{*}) = x^{*}\) for \(w \in \mathbf{W}_{\mathbf{X}}\) and \(x^{*} \in \mathbf{C}_{\mathbf{X}}\) .

PROPOSITION 5. Let \(X, X' \subset S\) and \(w, w' \in W\) . If \(w(C_X) \cap w'(C_{X'}) \neq \emptyset\) , then \(X = X'\) , \(wW_X = w'W_{X'}\) and \(w(C_X) = w'(C_{X'})\) .

We are reduced immediately to the case \(w' = 1\) . The proof is by induction on the length \(n\) of \(w\) . If \(n = 0\) , the assertion is clear. If \(l(w) > 0\) , there exists \(s \in S\) such that \(l(sw) = l(w) - 1\) and then (cf.~the end of no. 4) \(w(C) \subset s(A_s)\) , hence \(w(\overline{C}) \subset s(\overline{A_s})\) . Since \(\overline{C} \subset \overline{A_s}\) , it follows that

\[
\overline {{{\mathrm{C}}}} \cap w (\overline {{{\mathrm{C}}}}) \subset \mathrm{H} _ {s}.
\]

Then \(s(x^{*}) = x^{*}\) for all \(x^{*} \in \overline{\mathrm{C}} \cap w(\overline{\mathrm{C}})\) , and a fortiori for all

\[
x ^ {*} \in \mathrm{C} _ {\mathrm{X} ^ {\prime}} \cap w (\mathrm{C} _ {\mathrm{X}}).
\]

Consequently, the relation \(C_{X'} \cap w(C_X) \neq \varnothing\) implies on the one hand that \(C_{X'} \cap H_s \neq \varnothing\) , and hence that \(s \in X'\) , and on the other hand that \(C_{X'} \cap sw(C_X) \neq \varnothing\) . The induction hypothesis then implies that \(X = X'\) and \(swW_X = W_{X'} = W_X\) , so \(sw \in W_X\) and \(w \in W_X\) since \(s \in W_X\) . It follows that \(wW_X = W_{X'}\) and that \(w(C_X) = C_X = C_{X'}\) .

COROLLARY. Let X be a subset of S and \(x^{*}\) an element of \(C_{X}\) . The stabiliser of \(x^{*}\) in W is \(W_{X}\) .

Now let U be the union of the \(w(\overline{\mathbf{C}})\) for \(w \in W\) , and let F be the set of subsets of U of the form \(w(\mathrm{C}_{X})\) , with \(X \subset S\) and \(w \in W\) . By the above, F is a partition of U.

PROPOSITION 6. (i) The cone U is convex.

\begin{enumerate}
\def\labelenumi{(\roman{enumi})}
\setcounter{enumi}{1}
\item
  Every closed segment contained in U meets only finitely many elements of \(\mathfrak{F}\) .
\item
  The cone \(\overline{\mathbf{C}}\) is a fundamental domain for the action of W on U.
\end{enumerate}

To prove (iii), it is enough to show that, if \(x^{*}, y^{*} \in \overline{\mathbb{C}}\) and \(w \in \mathbb{W}\) are such that \(w(x^{*}) = y^{*}\) , then \(x^{*} = y^{*}\) . Now there exist two subsets \(\mathbf{X}\) and \(\mathbf{Y}\) of \(\mathbf{S}\) such that \(x^{*} \in \mathrm{C}_{\mathbf{X}}\) and \(y^{*} \in \mathrm{C}_{\mathbf{Y}}\) ; we have \(w(\mathrm{C}_{\mathbf{X}}) \cap \mathrm{C}_{\mathbf{Y}} \neq \emptyset\) , and Prop. 5 shows that \(\mathbf{X} = \mathbf{Y}\) and \(w \in \mathrm{W}_{\mathbf{X}}\) , which implies that \(x^{*} = y^{*}\) .

Now let \(x^{*}, y^{*} \in \mathrm{U}\) ; we shall show that the segment \([x^{*}y^{*}]\) is covered by finitely many elements of \(\mathfrak{F}\) , which will establish both (i) and (ii). By transforming \(x^{*}\) and \(y^{*}\) by the same element of \(\mathrm{W}\) , we can assume that \(x^{*} \in \overline{\mathrm{C}}\) . Let \(w \in \mathrm{W}\) be such that \(y^{*} \in w(\overline{\mathrm{C}})\) . We argue by induction on the length of \(w\) . For \(s \in \mathrm{S}\) , the relation \(w(\overline{\mathrm{C}}) \not\subset \overline{\mathrm{A}_{s}}\) is equivalent to \(w(\mathrm{C}) \notin \mathrm{A}_{s}\) and hence to \(l(sw) < l(w)\) (cf.~no. 4). Prop. 7 of no. 8 of Chap. IV, §1 now implies that there exist only finitely many \(s \in \mathrm{S}\) such that \(w(\overline{\mathrm{C}}) \not\subset \overline{\mathrm{A}_{s}}\) . The set \(T\) of \(s \in S\) such that \(\langle y^{*}, e_{s} \rangle < 0\) is thus finite. On the other hand, the intersection \(\overline{\mathrm{C}} \cap [x^{*}y^{*}]\) is a closed segment \([x^{*}z^{*}]\) . If \(z^{*} = y^{*}\) , that is if \(y^{*} \in \overline{\mathrm{C}}\) , there exist subsets \(X\) and \(Y\) of \(S\) such that \(x^{*} \in C_{X}\) and \(y^{*} \in C_{Y}\) . The open segment \(|x^{*}y^{*}|\) is then contained in \(C_{X \cap Y}\) , so \([x^{*}y^{*}] \subset C_{X} \cup C_{Y} \cup C_{X \cap Y}\) . If \(z^{*} \neq y^{*}\) , there exists an \(s \in T\) such that \(z^{*} \in H_{s}\) . Then \(w(C) \not\subset A_{s}\) and \(l(sw) < l(w)\) . The induction hypothesis thus implies that the segment \([z^{*}y^{*}] = s([z^{*}(s(y^{*}))])\) is covered by a finite number of elements of \(\mathfrak{F}\) , and hence so is

\[
\left[ x ^ {*} y ^ {*} \right] = \left[ x ^ {*} z ^ {*} \right] \cup \left[ z ^ {*} y ^ {*} \right]
\]

since \([x^{*}z^{*}]\subset \overline{\mathbf{C}}.\)

